\section{The case of unbounded overlap scale}\label{sec:unbounded}

The scale $W=W_2+A$ is finite and nonnegative for each configuration.
Every counterexample sequence therefore has a bounded-$W$ subsequence
or a subsequence on which $W\to\infty$. We treat the latter here;
Section~\ref{sec:bounded} treats the bounded alternative.
Suppose, then, that $W\to\infty$. We will derive a
contradiction from the variance estimate and the equation defining each
fiber. First we show that a positive fraction of $W$ comes from cells
whose uniform masses stay away from zero. These cells define finite
weighted trees. Transport on the trees will turn the fiber equation into
an impossible assertion about three independent uniform vertices.
Throughout this section, $o(W+1)=o(W)$.

\subsection{Extracting cells of positive mass}

We first claim that
\begin{equation}\label{eq:unb-overlap-scale}
 W_2\ge c_1 W
\end{equation}
for some fixed $c_1>0$ and all sufficiently large $n$. At a fixed
place and choice of grids, abbreviate $S=S_1+S_2$ and $L=L_1+L_2$.
Choose a fiber $e$ with its edge-induced weight and then choose $Y,Y'$
independently with law $\mu_e$. Equation~\eqref{eq:depth} and the
constant raw depth sum on the fiber give
\[
 S(Y)=M_1+M_2-S(x_e)-h_e+L(x_e,Y)
       -\Delta_1(x_e,Y)-\Delta_2(x_e,Y).
\]
All source--neighbor pairs are distinct. Thus the first-moment bound
from~\eqref{eq:grid-error}, the normalization of the archimedean
weights, and~\eqref{eq:q-overlap} imply
\[
 \avg{\E|S(Y)-S(Y')|}\le 2\avg{\E_q L(x,y)}+O(1)
 =O(W_2)+O(1).
\]
To compare this with $A$, we use the pointwise bound
\begin{equation}\label{eq:potential}
 S^*(x)\ge A-2W_2\qquad(n>2).
\end{equation}
Indeed, the average of $L^*(x,y)$ over the $n-1$ labels $y\ne x$ equals
\[
 W_2+\frac{n-2}{n-1}\bigl(S^*(x)-A\bigr),
\]
and is nonnegative.  Hence $\avg{\E S(Y)}\ge A-2W_2$.
Consequently
\begin{equation}\label{eq:unb-min-lower}
 \avg{\E\min\{S(Y),S(Y')\}}\ge A-O(W_2)-O(1).
\end{equation}

For two nonnegative lists $(a_1,a_2)$ and $(b_1,b_2)$,
\[
 \min\{a_1+a_2,b_1+b_2\}\le\sum_{i,j=1}^2\min\{a_i,b_j\}.
\]
Indeed, if the first sum is smaller, then
$\sum_j\min\{a_i,b_j\}\ge a_i$ for each $i$. Moreover, nesting of the
giants gives, for almost every $u>0$,
\[
 \{z:S_i(z)>u\}=P\setminus G_i(M_i-u).
\]
A vertex which leaves the giant at level $r_z<M_i$ has
$S_i(z)=M_i-r_z$; one which never leaves has $S_i(z)=0$. This proves
the displayed identity, with only finitely many possible threshold ties.
Conditional independence of $Y,Y'$ and $ab\le(a^2+b^2)/2$ now give
\begin{equation}\label{eq:unb-min-upper}
 \avg{\E\min\{S(Y),S(Y')\}}
 \le 2\avg{\sum_{i=1}^2\E_e\int_{r<M_i}
          \mu_e(P\setminus G_i(r))^2\,dr}.
\end{equation}
To bound the right side, write $p=\mu_e(C)$ and $a=\alpha(C)$ for a
giant-complement occurrence. If $|C|>n^{1-c}$, use
$p^2\le2(p-a)^2+2a^2$, the variance estimate~\eqref{eq:variance}, and
$a^2\le2\pi(C^2)$. For smaller complements with $p\ge2/k_e$, use
the second term of~\eqref{eq:variance}. The remaining terms satisfy
$p^2\le2n^{-\gamma}p$. Since $\E_e p=q_Y(C)\le B_0a$, their integral
is $O(n^{-\gamma}A)$. In particular, this last error is charged to
$A$, the finite integral of the complement masses. We obtain
\[
 \avg{\E\min\{S(Y),S(Y')\}}
 \le O(W_2)+o(W)+O(n^{-\gamma}A).
\]
Together with~\eqref{eq:unb-min-lower}, this proves
$A\le O(W_2)+o(W)+O(1)$, and hence~\eqref{eq:unb-overlap-scale}.
Equation~\eqref{eq:q-overlap} also gives $\I(q(C^2))\gg W$.

We next restrict the masses of the cells we use. For a fixed
$0<\delta<1/4$, split the contribution of $a\le\delta$ according to
$p<2\delta$ or $p\ge2\delta$. In the latter case
$p\le p^2/(2\delta)$ and $p-a\ge p/2$. On large cells this is
controlled by the first variance term; on small cells,
$p\ge2\delta\ge2/k_e$ eventually, so the second variance term applies.
Therefore
\[
 \I\bigl(\1_{\{a\le\delta\}}q(C^2)\bigr)
 \le T(2\delta)+o_\delta(W)
 \le O(\delta^{1/3}W)+o_\delta(W).
\]
For the last inequality apply Lemma~\ref{lem:small-fiber} with
$\lambda=2\delta$, $\tau=\lambda^{1/3}$ and any fixed $H\ge1$.
The terms depending only on $\delta,H$ are $o_\delta(W)$.
The implicit constant in $O(\delta^{1/3}W)$ is independent of $\delta$.
We may thus choose $\delta$ sufficiently small once and for all, and
then let $n\to\infty$, to obtain
\begin{equation}\label{eq:cut-mass}
 \I(\1_{\{a>\delta\}})\gg W,
 \qquad
 \I(\1_{\{a\ge\delta/4\}})\ll_\delta W.
\end{equation}
The first inequality uses $q(C^2)\le1$. For the second, when $n$ is
large, $a\ge\delta/4$ implies
$\pi(C^2)\ge a^2/2\ge\delta^2/32$.
All cutoffs below use this fixed $\delta$; in particular, a bound
$O_\delta(1)$ is $o(W)$ along the present subsequence.

Choose a Lipschitz function $w:[0,1]\to[0,1]$ satisfying
\[
 w(u)=w(1-u),\qquad
 w=0\text{ on }[0,\delta/2]\cup[1-\delta/2,1],\qquad
 w=1\text{ on }[\delta,1-\delta].
\]
For the original partitions, define
\[
 \ell_j=\int_{\R}\sum_{B\in\Pi_j(r)}w(\alpha(B))\,dr.
\]
Replacing a giant by its complement preserves every summand by
symmetry, with occurrences still counted separately. Thus
\begin{equation}\label{eq:unb-total-length}
 W\ll\avg{\ell_1+\ell_2}\ll_\delta W.
\end{equation}

For either family of measures in Lemma~\ref{lem:variance}, we shall use
\begin{equation}\label{eq:cut-approx}
 \avg{\E_{\mathcal F}\sum_{j=1}^2\int_{\R}
 \sum_{B\in\Pi_j(r)}
 \left[|w(\nu(B))-w(\alpha(B))|
       +w(\alpha(B))|\nu(B)-\alpha(B)|\right]dr}=o(W).
\end{equation}
The first term lets us replace cut weights defined by the fiber laws
with those defined by the uniform law.  The second controls the weighted
mass imbalance that will appear in the transport comparison.
Here the family of incoming laws retains only targets with at least
$n^\gamma$ incoming neighbors, as in Section~\ref{sec:variance}.
To prove~\eqref{eq:cut-approx}, note first that its integrand is
unchanged under $(p,a)\mapsto(1-p,1-a)$. Work therefore with the
transformed cells. On $a\ge\delta/4$, all cells are large for large
$n$; the Lipschitz bound, Cauchy--Schwarz,
\eqref{eq:cut-mass}, and~\eqref{eq:variance} give $o(W)$.
On $a<\delta/4$, the second summand vanishes, and the first can be
nonzero only if $p>\delta/2$. For large cells,
$1\le16\delta^{-2}(p-a)^2$; for small cells,
$1\le4\delta^{-2}p^2$ and $p\ge2/k$ eventually. The same variance
estimate again gives $o(W)$.

\subsection{Distances obtained from the hierarchy}

We have retained a positive total amount of cells, and the masses of
those cells are close to their $\alpha$-masses in the fiber average. We now use these cells as cuts defining a distance. For
$\nu=\alpha$ or $\mu_e$, put
\[
 D_j^\nu(a_0,a_1)=\int_{\R}\sum_{B\in\Pi_j(r)}
 |\1_B(a_0)-\1_B(a_1)|\,w(\nu(B))\,dr.
\]
We also allow a formal endpoint $\infty$, setting $\1_B(\infty)=0$
for every cell. These quantities are finite for all sufficiently large
$n$. On the initial ray the only cell is $P$, of weight $w(1)=0$;
on the final ray all cells are singletons, of $\nu$-mass either
$0$, $1/n$, or $1/k_e$, so their weights vanish too.

For distinct $a_0,a_1\in P$ we have
\begin{equation}\label{eq:projection}
 D_j^\nu(a_0,a_1)=\int_{\R}
 w\bigl(\nu\{z:d_j(a_1,z)-d_j(a_0,z)\ge u\}\bigr)\,du.
\end{equation}
Indeed, set $d_0=d_j(a_0,a_1)$. Nestedness implies that if
$d_j(a_1,z)>d_0$, then $d_j(a_0,z)=d_0$, and conversely with the
indices exchanged. Otherwise the two depths to $z$ agree.
For $u>0$ the event in~\eqref{eq:projection} is the $a_1$-cell at
level $d_0+u$. For $u<0$ its complement is the $a_0$-cell at level
$d_0-u$, up to strict versus weak threshold ties. There are only
finitely many finite depths, so those ties have zero Lebesgue measure.
Splitting the integral at zero and using $w(t)=w(1-t)$ gives exactly
the two cell contributions above their common depth $d_0$.
The depth difference is $+\infty$ at $z=a_1$ and $-\infty$ at
$z=a_0$; no difference $\infty-\infty$ occurs. Directly from the
cell definition, we also have
\begin{equation}\label{eq:unb-root-projection}
 D_j^\nu(\infty,y)=\int_{\R}
 w\bigl(\nu\{z:d_j(y,z)\ge u\}\bigr)\,du.
\end{equation}

For an edge $(x,y)$ in fiber $e$, put $j'=3-j$. The fiber equation
implies, for $z\in F_e$,
\[
 Z_{j'}(z)=Z_{j'}(x)+\frac{t_e}{Z_j(z)-Z_j(x)}.
\]
Consequently, for $z\ne y$ in this fiber,
\[
 Z_{j'}(z)-Z_{j'}(y)
 =\frac{t_e\bigl(Z_j(y)-Z_j(z)\bigr)}
 {(Z_j(z)-Z_j(x))(Z_j(y)-Z_j(x))},
\]
and hence
\begin{equation}\label{eq:unb-raw-fiber}
 i_{j'}(y,z)=i_j(y,z)-i_j(x,z)+h_e-i_j(x,y).
\end{equation}
This identity converts one of the projected distances into a root
distance in the other coordinate. Here are the details, including the
diagonal atom. Under $z\sim\mu_e$, compare
\[
 U(z)=d_{j'}(y,z),\qquad
 V(z)=d_j(y,z)-d_j(x,z)+h_e-i_j(x,y).
\]
For $z\ne y$, Equation~\eqref{eq:unb-raw-fiber} gives
\[
 |U(z)-V(z)|\le
 |\Delta_{j'}(y,z)|+|\Delta_j(y,z)|+|\Delta_j(x,z)|.
\]
At $z=y$, both variables are $+\infty$, so their indicators at every
finite threshold agree; $z=x$ has zero $\mu_e$-mass. Tonelli's theorem
therefore gives the precise bound
\[
 \int_{\R}|\Pr(U\ge u)-\Pr(V\ge u)|\,du
 \le \sum_{z\in F_e\setminus\{y\}}\mu_e(z)|U(z)-V(z)|.
\]
Translation of the threshold by $h_e-i_j(x,y)$ does not change the
integral of the weighted tail. Apply the Lipschitz bound for $w$,
\eqref{eq:projection}, and~\eqref{eq:unb-root-projection}. After
averaging over $e$, $y\sim\mu_e$, grids and places, we obtain
\[
 \avg{\E_q
 |D_j^{\mu_e}(x,y)-D_{j'}^{\mu_e}(\infty,y)|}=O_\delta(1).
\]
The error estimate is valid for these sampling laws because
\eqref{eq:grid-error} bounds the grid first moment uniformly for
every distinct pair. All vertex laws are independent of the grids;
nonarchimedean errors vanish and the archimedean weights sum to one.

Changing $\mu_e$ to $\alpha$ changes any one of its cut distances by
at most $\int\sum_B|w(\mu_e(B))-w(\alpha(B))|\,dr$, independently
of the endpoints. Thus~\eqref{eq:cut-approx}, with
$D_j=D_j^\alpha$, yields
\begin{equation}\label{eq:fiber-tree}
 \avg{\E_q|D_j(x,y)-D_{j'}(\infty,y)|}=o(W).
\end{equation}
The remaining task is to replace the edge law in this equation by
independent uniform vertices and then exploit the geometry of the
cut distances.

\subsection{Transport to independent uniform vertices}

Each $D_j$ is the path distance on a finite weighted tree, allowing
zero-length edges. To see this, take the distinct cells of the nested
partitions as vertices, root the inclusion tree at $P$, and connect
each proper cell $B$ to its least strict supercell. Give this edge
length
\[
 l_j(B)=\int_{\R}\1_{\{B\in\Pi_j(r)\}}w(\alpha(B))\,dr.
\]
The points of $P$ are singleton leaves. An edge is crossed by their
connecting path exactly when its cell separates them, proving the
cut formula for $D_j$. The total edge length is exactly $\ell_j$.
The formal endpoint is represented by the root, since the root-cell
weight is zero.

For probabilities $\rho,\alpha$ on the leaves, let
$\mathsf T_j(\rho,\alpha)$ be the least expected $D_j$-distance among
couplings of the two laws. The tree transport formula is
\begin{equation}\label{eq:unb-transport}
 \mathsf T_j(\rho,\alpha)
 =\sum_{B\ne P}l_j(B)|\rho(B)-\alpha(B)|
 =\int_{\R}\sum_{B\in\Pi_j(r)}
       w(\alpha(B))|\rho(B)-\alpha(B)|\,dr.
\end{equation}
This standard formula is the finite-tree form of
\cite[Section~2, Equation~(5)]{TreeTransport}. It has a
short direct proof. Across the edge above $B$, any coupling must move
at least $|\rho(B)-\alpha(B)|$ units of mass. Conversely, first match
common mass at each leaf. Work upward, matching surplus mass from child
subtrees to deficits in other child subtrees, and send only the
remaining imbalance toward the parent. Across each edge the total
mass moved is precisely its absolute imbalance; the root imbalance
is zero. Decomposing these finitely many transfers into leaf-to-leaf
paths gives a coupling attaining the displayed cost. Zero-length
edges do not affect the argument.

By $q_Y=\E_e\mu_e$, convexity and~\eqref{eq:cut-approx} give
\begin{equation}\label{eq:unb-two-costs}
 \avg{\mathsf T_j(q_Y,\alpha)}=o(W),\qquad
 \avg{\E_{q_Y}\mathsf T_j(\rho_y,\alpha)}=o(W).
\end{equation}
For the second estimate, first apply~\eqref{eq:cut-approx} to targets
with at least $n^\gamma$ incoming neighbors. The omitted $q_Y$-mass
is $O(n^{\gamma-1})=o(1)$, independent of the hierarchy. Every
transport cost is at most $\ell_j$, so its omitted contribution is
$o(1)\avg{\ell_j}=o(W)$ by~\eqref{eq:unb-total-length}.

The conditional source law in Equation~\eqref{eq:fiber-tree}, given
the target $y$, is $\rho_y$. We keep this conditioning while coupling
the sources to the uniform law.
Fix a hierarchy for a moment and put
$F(x,x',y)=|D_j(x,y)-D_j(x',y)|$. If $Y\sim q_Y$ and $X,X'$ are
independent with conditional law $\rho_Y$, then
\[
 \E F(X,X',Y)
 \le2\E_q|D_j(x,y)-D_{j'}(\infty,y)|.
\]
For each $y$, use two independent copies of an optimal coupling of $\rho_y$
to $\alpha$, replacing $X,X'$ by $U,U'$. Conditional on $y$, the
new points are independent with law $\alpha$; since that law does
not depend on $y$, they are also independent of $Y$. The triangle
inequality bounds the change in $F$ by the sum of the two transport
distances. Next optimally couple $Y\sim q_Y$ to $V\sim\alpha$, independently
of $U,U'$. Moving the last argument of $F$ costs at most twice its
transport distance. We have proved the pointwise inequality
\begin{align*}
 \E_{\alpha^{\otimes3}}F
 &\le 2\E_q|D_j(x,y)-D_{j'}(\infty,y)|\\
 &\quad+2\E_{q_Y}\mathsf T_j(\rho_y,\alpha)
       +2\mathsf T_j(q_Y,\alpha).
\end{align*}
Its right side consists of explicit finite cut sums, so no selection
of couplings across places or random grids is required. Integrating
and using~\eqref{eq:fiber-tree} and~\eqref{eq:unb-two-costs} gives
\begin{equation}\label{eq:iid-tree}
 \avg{\E_{x,x',y\ \mathrm{iid}\ \alpha}
 |D_j(x,y)-D_j(x',y)|}=o(W).
\end{equation}

\subsection{A lower bound for tree distance fluctuations}

We finish with an elementary fact whose constants do not depend on
the shape of the tree.

\begin{lemma}\label{lem:unb-tree-fluctuation}
Let a finite tree have nonnegative edge lengths of total sum $L$,
with path pseudometric $d$, and let $\alpha$ be a probability on its
vertices. Suppose every edge of
positive length separates two sets each having $\alpha$-mass at least
$\delta/2$, where $0<\delta<1$. If $X,X',Y$ are independent with
law $\alpha$, then
\[
 \E|d(X,Y)-d(X',Y)|\ge
 \frac{\delta^2}{32(\lceil16/\delta\rceil+1)}L.
\]
\end{lemma}

\begin{proof}
There is nothing to prove if $L=0$. For every fixed $y$, the expected
distance from $X$ to $y$ is at least $\delta L/2$: each positive
edge separates $y$ from mass at least $\delta/2$. Since all distances
are at most $L$,
\[
 \Pr\{d(X,y)\ge\delta L/4\}\ge\delta/4.
\]
Traverse the tree in a closed walk crossing each edge twice, of total
length $2L$. Partition this walk into at most
$M=\lceil16/\delta\rceil+1$ intervals of length at most $\delta L/8$,
and assign each vertex to an interval containing one chosen visit.
The resulting classes have diameter at most $\delta L/8$. Thus
\[
 \Pr\{d(X',Y)\le\delta L/8\}
 \ge\sum_{r=1}^M\alpha(C_r)^2\ge1/M.
\]
Conditional on $Y$, the far event for $X$ and the near event for $X'$
are independent. The far probability is uniformly at least $\delta/4$,
so both events occur with probability at least $\delta/(4M)$.
On their intersection the absolute difference of the distances is
at least $\delta L/8$, proving the claim.
\end{proof}

In the tree for $D_j$, every positive edge has both side masses greater
than $\delta/2$, by the support condition on $w$. Apply
Lemma~\ref{lem:unb-tree-fluctuation} with $L=\ell_j$, sum over
$j=1,2$, and average over places and grids. Equations
\eqref{eq:unb-total-length} and~\eqref{eq:iid-tree} give
$W\ll o(W)$, a contradiction. This excludes every subsequence on
which $W\to\infty$.
